% from main.tex
%%%%%%%%%%  scan idx 161  |  folio 148  %%%%%%%%%%
\emph{Preuve.} La première assertion résulte trivialement des relations \(\cl(L[1])=-\cl(L)\), \(\cl(L^n)=n\cl(L)\) pour \(L\in\Ob\Parf(S)\). Prouvons la seconde. On remarque d'abord que \(i\) se factorise de manière évidente à travers \(k(\mathcal C_0)\). Comme \(\rg\, i\) est l'identité, on est ramené (d'après (6.9)) à prouver que l'homomorphisme \(i\, \rg:k(\mathcal C_0)\longrightarrow k(\mathcal C_0)\) induit un isomorphisme sur les faisceaux associés, i.e. que, pour \(L\in\Ob\mathcal C_0\), \(i\, \rg(L)\) est localement isomorphe à \(L\), ce qui est immédiat.

On en déduit, d'après (6.6.b)),

\textbf{Corollaire 6.15.} Sous les hypothèses de (6.13), pour tout faisceau abélien \(F\) sur \(S\), le groupe des \(S\)-fonctions additives de \(\Parf(S)\) dans \(F\) s'identifie canoniquement à \(H^0(S,F)\) (\(S\) = objet final de \(S\)), la fonction rang correspondant à \(F=\mathbb Z_S\) et à la section unité de \(\mathbb Z_S\).

\section*{7. Dualité des complexes parfaits}

\textbf{7.0.} Dans ce numéro, \((S,A)\) désigne un topos annelé en anneaux commutatifs, on note \(S\) l'objet final de \(S\). On se propose de généraliser aux complexes parfaits les notions de dualité bien connues pour les faisceaux localement libres de type fini.

\textbf{Proposition 7.1.} Soient \(E\in\Ob\Dpx(S)_{\mathrm{parf}}\), \(F\in\Ob\Dpx^+(S)\).

\textbf{a)} Si \(F\) est à cohomologie localement bornée, il en est de même de \(\RHom(E,F)\). Si \(F\in\Ob\Dpx(S)\) est pseudo-cohérent (resp. localement de tor-dimension finie, resp. parfait), il en est de même de \(\RHom(E,F)\).

\textbf{b)}
